% SHORT IEEE 6-page target. Historical results unverified, not for final submission.
% !TeX program = xelatex
\documentclass[conference]{IEEEtran}

\usepackage{fontspec}
\setmainfont{TeX Gyre Termes}
\setsansfont{TeX Gyre Heros}
\usepackage{graphicx}
\usepackage{amsmath}
\usepackage{amssymb}
\usepackage{array}
\usepackage{booktabs}
\usepackage{multirow}
\usepackage{tikz}
\usepackage{pgfplots}
\usepackage{pgf-pie}
\pgfplotsset{compat=1.17}
\usetikzlibrary{arrows.meta, positioning, shapes.geometric}

\begin{document}

\title{CropState-Vision: Reliable Six-Stage Rice Growth
Classification and Confidence-Aware Stage-Gated Retrieval from Field Images}

\author{\IEEEauthorblockN{Lê Bảo Nhi \qquad Phan Nguyễn Quốc Minh \qquad Lê Khánh Nhi}
\IEEEauthorblockA{FPT University, Da Nang, Viet Nam\\
\{lebaonhi0805, minhpnq1807, nhilekhanh29\}@gmail.com}}

\maketitle

\begin{abstract}
Agricultural recommendations depend on crop development, so a decision-support system must first recognize the growth stage and then retrieve stage-appropriate guidance. This paper presents CropState-Vision, which classifies field images into six BBCH-derived rice macro-stages and uses the full stage distribution to gate evidence retrieval. Beyond raw accuracy, our emphasis is evaluation reliability and retrieval safety. On a 2{,}900-image benchmark derived from RiceSEG, we first audit the splits for near-duplicate leakage and detect no cross-split near-duplicates at perceptual-hash thresholds up to 10, reducing concern that the reported scores are inflated by near-duplicate memorization. Using a fixed-split, three-seed protocol we compare three training objectives for a ResNet-18 classifier: cross-entropy reaches 0.800 accuracy, an ordinal-aware loss leaves overall accuracy unchanged but stabilizes recall on the under-represented reproductive stage (identical across the three seeds, $\sigma=0.000$ versus $0.052$ for cross-entropy), and a focal loss degrades calibration and is kept only as a negative ablation. We recast confidence handling as selective prediction: abstaining on the least-confident ${\sim}46\%$ of cases roughly halves the error rate ($0.205\to0.099$). Finally, on 420 automatically-derived retrieval scenarios, confidence-adaptive soft gating remains competitive in ranking quality while reducing stage-incompatible retrievals by approximately 75\% relative to query expansion; an error-stratified analysis, with paired image-clustered bootstrap tests, shows this ranking-versus-safety trade-off holds when the prediction is correct or off by one adjacent stage, while hard filtering stays safest under rare non-adjacent errors. Because these relevance labels are auto-derived from the same stage-compatibility metadata the re-ranker uses, we interpret the retrieval results as mechanism validation---consistency with stage-compatibility metadata---rather than evidence of expert-judged agronomic effectiveness. CropState-Vision claims no new architecture or dataset; the contribution is a reliable, safety-oriented evaluation of the full classification-to-retrieval pipeline.

\begin{IEEEkeywords}
rice growth stage, BBCH, image classification, ordinal loss, calibration,
selective prediction, stage-aware retrieval, data leakage
\end{IEEEkeywords}
\end{abstract}

\section{Introduction}

Agricultural management follows the crop, not the calendar. Water management, nutrient application, disease surveillance, and harvest preparation all differ across establishment, tillering, reproductive development, grain filling, and ripening, so a digital advisor needs to know which stage a field is in before it can retrieve safe, timely guidance.

Field images are a natural input, because phenological development shows up visually through leaf development, tiller formation, stem elongation, panicle emergence, grain formation, and ripening. In practice, stage recognition is hardest exactly where it matters: near stage boundaries, and under everyday variation in viewpoint, lighting, plant density, background, and camera distance. A model can also be confidently wrong, which is a particular risk for a class with few training examples. Converting an uncertain visual prediction into a single stage label and applying a hard metadata filter can then amplify the error by discarding all evidence associated with the true stage.

CropState-Vision represents crop development as a distribution over stages rather than a single static label. Image classification supplies the visual evidence; the distribution becomes a stage belief that controls how strongly the retriever favors stage-compatible evidence. When the prediction is concentrated and calibrated, stage compatibility receives more weight; when it is diffuse, the system relaxes stage conditioning to preserve recall. After an image is uploaded, the system retrieves evidence for a fixed set of care topics rather than answering free-form questions.

This paper asks two questions. First, how reliably can six rice macro-stages be recognized from field images, once the benchmark is audited for leakage and the classifier is evaluated across seeds rather than on a single lucky run? Second, is confidence-aware soft gating more robust than top-1 hard filtering when the image-based prediction is uncertain or wrong? The contributions are:
\begin{enumerate}
    \item a near-duplicate-audited benchmark: an average-hash (aHash) near-duplicate sweep over 2{,}900 images detects no cross-split pairs at Hamming thresholds up to 10 under an aHash-grouped split; this property is partly enforced by construction and does not establish field-level independence;
    \item a reliable stage-classification evaluation using a fixed-split, multi-seed protocol that separates initialization variance from split variance, showing that an ordinal-aware loss stabilizes reproductive-stage recall without changing accuracy, plus a selective-prediction analysis showing the error rate roughly halves when the model abstains on its least-confident cases;
    \item a confidence-aware soft-gating retriever that uses the full visual distribution instead of only the top-1 stage, and an end-to-end evaluation over 420 auto-derived scenarios---with paired, image-clustered bootstrap tests---showing it stays competitive in ranking quality while cutting stage-incompatible retrievals by approximately 75\% relative to the strongest baseline;
    \item an explicit, honest account of the limitations---an underpowered reproductive subset reported descriptively rather than pooled, retrieval labels auto-derived from the same metadata the re-ranker uses (mechanism validation, not expert-judged effectiveness), and simulated temporal fusion---that bound how far these results can be pushed.
\end{enumerate}

\section{Related Work}
The BBCH scale describes observable phenological transitions using hierarchical stage codes \cite{meier2001bbch,lancashire1991uniform}; our six categories aggregate these codes and do not replace expert staging. Image-based rice growth work has considered seedling classification and reproductive-stage panicle detection \cite{tan2022machine,tan2023ricepanicle,zheng2025android}, whereas our task covers establishment through ripening using whole-field images. We adopt an ImageNet-pretrained ResNet-18 \cite{he2016deep,deng2009imagenet} and evaluate ordinal errors, class imbalance and temperature scaling \cite{gutierrez2016ordinal,lin2017focal,guo2017calibration}. Retrieval combines BM25, dense embeddings and reciprocal rank fusion \cite{robertson2009bm25,karpukhin2020dpr,cormack2009rrf}; our focus is how uncertainty in visual stage beliefs affects stage-conditioned reranking and metadata-defined incompatibility, not a new image architecture or retrieval foundation model.

\section{Task Formulation}

\subsection{Research Questions and Hypotheses}

RQ1 tests whether pretrained classification surpasses a majority baseline; RQ2 evaluates whether confidence-aware soft gating improves the ranking-versus-incompatibility trade-off relative to hard filtering; RQ3 examines ordinal errors and rare-class limitations.


% Condensed for 6-page IEEE draft: tab:rq retained in full manuscript / supplement.


\subsection{Six Rice Macro-Stages and Stage Belief}

Let $\Phi = \langle \phi_1,\dots,\phi_6 \rangle$ be the ordered stage set in Table~\ref{tab:stages}. For image $I$, a model $f_\theta$ produces logits $o = f_\theta(I)$, and the stage distribution is
\begin{equation}
p(\phi_i \mid I) = \mathrm{softmax}(o/T)_i,
\label{eq:softmax}
\end{equation}
where $T > 0$ is a temperature estimated from validation data when calibration is applied. The predicted class is $\hat{\phi} = \arg\max_i p(\phi_i \mid I)$, but CropState-Vision keeps the full distribution $B=p(\cdot\mid I)$ as the stage belief. Belief concentration is summarized by normalized entropy $\Gamma = 1 - H(B)/\log 6$ and the top-two margin $M = B^{(1)} - B^{(2)}$, combined into a confidence proxy $\kappa = \eta_1\Gamma + \eta_2 M$ with $\eta_1+\eta_2=1$. This confidence is a routing signal, not proof that the predicted stage is correct.

\begin{table}[t]
\centering
\caption{Rice macro-stages used in the study.}
\label{tab:stages}
\begin{tabular}{lcl}
\toprule
\textbf{Group} & \textbf{BBCH} & \textbf{Label} \\
\midrule
Germination and leaf development & 00--19 & Establishment \\
Tillering & 20--29 & Tillering \\
Stem elongation and booting & 30--49 & Stem/booting \\
Heading and flowering & 50--69 & Reproductive \\
Grain development & 70--79 & Grain filling \\
Ripening & 80--89 & Ripening \\
\bottomrule
\end{tabular}
\end{table}

\subsection{Confidence-Aware Stage-Gated Retrieval}

Each knowledge chunk $x$ carries a stage-compatibility vector $C(x,\phi_i)\in[0,1]$, where one denotes direct applicability, an intermediate value adjacent or multi-stage applicability, and zero incompatibility or contraindication. For a fixed care topic, BM25 and dense retrieval produce candidate lists whose RRF score is min--max normalized to $\overline{S}_{\mathrm{base}}(x)$. The expected stage compatibility is
\begin{equation}
S_{\mathrm{stage}}(x,B)=\sum_{i=1}^{6}B(\phi_i)\,C(x,\phi_i).
\end{equation}
The adaptive stage weight rises with confidence, $\beta=\beta_{\min}+(\beta_{\max}-\beta_{\min})\kappa$, and the final score is
\begin{equation}
S(x)=(1-\beta)\,\overline{S}_{\mathrm{base}}(x)+\beta\,S_{\mathrm{stage}}(x,B).
\end{equation}
Unlike hard filtering, this never excludes a document solely because it is incompatible with the top-1 stage; hard filtering is retained as a comparison baseline. Stage safety is measured by the Stage-Inappropriate Retrieval Rate, $\mathrm{SIRR}@k=\frac{1}{k}\sum_{x\in\mathrm{TopK}}\mathbb{I}[C(x,\phi^*)=0]$, where $\phi^*$ is the ground-truth stage.

\section{Materials and Methods}

\subsection{Dataset and Leakage Audit}

The evaluation uses a 2{,}900-image benchmark derived from RiceSEG, split by group (by parent image, so overlapping crops of one photograph stay on the same side) into 2{,}400 training, 237 validation, and 263 test images (seed 42). The split is held \emph{fixed} across all model-seed reruns, so the multi-seed analysis measures training variability conditional on this fixed partition; it does not estimate variance across dataset splits. Class counts are imbalanced across all splits (test: establishment 54, tillering 88, stem/booting 64, reproductive 24, ripening 23, grain filling 10), most severely for grain filling and the reproductive stage.

RiceSEG ships per-pixel segmentation masks but no growth-stage labels; we derive the six macro-stages from mask morphology (canopy cover, panicle fraction, senescence fraction) using BBCH-motivated thresholds. Concretely, images with an emerged panicle (BBCH~$\ge$~50) are assigned to reproductive, grain filling, or ripening by their senescence fraction, and pre-panicle images (BBCH~00--49) are assigned to establishment, tillering, or stem/booting by canopy cover. Because the pre-panicle split rests on canopy cover---a weak proxy for the tiller-count and stem-elongation criteria BBCH actually uses---39\% of images (1{,}134/2{,}900) are auto-flagged for expert review, concentrated in the under-represented classes (e.g.\ reproductive 222/288, grain filling 84/110). These labels have not been independently double-annotated; we state this weak-supervision limitation explicitly (Section~\ref{sec:limitations}).

To reduce the concern that scores are inflated by near-duplicate memorization, we audit the splits with a 64-bit average-hash (aHash) near-duplicate detector: images are converted to grayscale, resized to $8\times8$ pixels, thresholded at the mean grayscale intensity, and paired by Hamming distance. Connected components, together with common parent-image IDs, are kept in the same split \emph{before} the audit. Sweeping the Hamming threshold (aHash threshold sweep) and manually inspecting matches at thresholds 10 and 12, we find that at the reporting threshold (Hamming $\le 10$) the corpus contains 661 near-duplicate edges and 122 multi-image clusters, and \emph{no} such cluster straddles splits, as expected from the hash-aware split construction; cross-split matches appear only at looser thresholds ($\ge 12$) that already begin to link visually distinct images. We therefore report the splits as free of \emph{detected cross-split near-duplicates} rather than leak-free: this audit only addresses near-duplicate memorization and does not rule out field-, session-, season-, or background-level leakage. It is complementary to, not a replacement for, field- or season-grouped splitting, which remains a data-collection item (Section~\ref{sec:limitations}) because RiceSEG does not carry field/session identifiers.


% Condensed for 6-page IEEE draft: tab:leakage retained in full manuscript / supplement.


\subsection{Knowledge Base and Retrieval Labels}

The retrieval evaluation runs over the non-restricted subset of an English-language corpus of machine-curated chunks built from the publicly available IRRI Rice Knowledge Bank, each tagged with a care topic and a stage-compatibility vector; the corpus is used for research under the source's terms, and derived chunk tags are released with the code (Section~\ref{sec:availability}). For each (test image, care topic) pair we form a scenario using the image's six-class belief; chunks whose stage-compatibility with the ground-truth stage is high are labeled relevant and stage-incompatible chunks are labeled inappropriate. Keeping only (topic, stage) pairs the corpus can answer (at least three relevant and three incompatible chunks) yields 420 scenarios. Crucially, these relevance labels are \emph{auto-derived from the same stage-compatibility metadata that the proposed re-ranker consumes}, so the evaluation is a \emph{mechanism validation}---it measures consistency with the designed metadata---and cannot, on its own, establish agronomic relevance or recommendation effectiveness. We accordingly avoid expert-effectiveness language and treat every retrieval number as a preliminary automated stage-compatibility result pending an independent, expert-annotated relevance set.

\subsection{Vision Models and Loss Functions}

The backbone is fixed to a pretrained, fine-tuned ResNet-18; we vary the training objective, since the interest is stage structure and reliability rather than architecture search. Against a majority-class reference we compare standard cross-entropy, an ordinal-aware loss that adds a squared penalty on the expected stage index, $\mathcal{L}=\mathrm{CE}+\lambda(\mathbb{E}[\phi]-\phi^*)^2$, and a focal loss that down-weights easy examples. All configurations share the fixed group split, class-balanced sampling, and evaluation protocol, and each is trained with three seeds (42, 7, 123) so reported figures are mean $\pm$ standard deviation. Temperature scaling is fitted after training on validation logits only.

\subsection{Retrieval Baselines}

The five comparison methods are summarized below.


The comparison uses B0 (ungated hybrid BM25+dense RRF), B1 (top-1 stage query expansion), B2 (top-1 hard filtering), B3 (fixed-soft stage reranking), and P (confidence-adaptive belief-based reranking).

All methods use the same corpus, chunking, embedding model, candidate depth, fixed topics, and top-$k$; only the stage-conditioning mechanism differs.

\subsection{Evaluation Metrics and Statistics}

Vision metrics are accuracy, macro-F1, per-class recall, ECE, Brier score, and Mean Absolute Stage Distance, $\mathrm{MASD}=\frac{1}{N}\sum_j|\mathrm{index}(\hat{\phi}_j)-\mathrm{index}(\phi_j^*)|$, which separates an adjacent mistake from one several stages off. Deployment behavior is summarized by selective prediction: after temperature scaling on the validation split, we rank test images by the confidence proxy and, for each coverage level, report the selective risk ($1-$ selective accuracy). We summarize the curve by the normalized area under the risk--coverage curve (AURC), the trapezoidal integral of selective risk over coverage in $[0,1]$, where lower is better. Retrieval metrics are Precision@$k$, Recall@$k$, nDCG@$k$, and SIRR@$k$ ($k=5$). Because methods are evaluated on the same items, comparisons are paired: for vision we report bootstrap 95\% confidence intervals \cite{efron1993bootstrap} and McNemar's test \cite{mcnemar1947note} for paired binary outcomes. The reproductive-recall behavior across the three seeds reuses the same 24 test images, so those repeated predictions are \emph{not} independent observations; we therefore report them descriptively (per-seed recall and its spread) and do not pool them for a single significance test. For retrieval, because one test image spawns several topic scenarios, we use a paired bootstrap that \emph{resamples whole images} (clusters of scenarios) and apply a Holm correction across the family of method comparisons, so the reported confidence intervals respect the non-independence of same-image scenarios.

% CAMERA-READY REPRODUCTION HOLD: Tables retain legacy reported numbers.\n% Re-run 420 scenarios with frozen corpus/logits after calibration, source-weight,\n% and candidate-universe corrections (see docs/RIVF26_REVIEWER_REVISION.md).\n\section{Results}

\subsection{Stage Classification (RQ1)}

All three losses beat the majority-class reference (accuracy 0.335) by a wide margin (Table~\ref{tab:results}). Among the losses, differences in overall accuracy, macro-F1, and MASD are within seed noise: the ordinal loss holds a razor-thin, non-meaningful lead on mean accuracy (0.801 vs 0.800), macro-F1 (0.796 vs 0.794), and MASD (0.233 vs 0.238), while cross-entropy has the lowest ECE. A single seed can even invert the accuracy ranking---on seed 42 the ordinal model reaches 0.814 versus cross-entropy's 0.795, whereas cross-entropy leads on the highest-accuracy seed---which is why we report three-seed means on a fixed split rather than a single lucky run.

\begin{table}[t]
\centering
\caption{ResNet-18 on the 2{,}900-image benchmark (2{,}400 train / 237 validation / 263 test; fixed group split). Each learned row is mean $\pm$ standard deviation over three seeds (42, 7, 123) sharing the \emph{same} split, so the spread reflects initialization variance only. Best value per column in bold.}
\label{tab:results}
\small
\begin{tabular}{lccccc}
\toprule
\textbf{Model / loss} & \textbf{Acc.} & \textbf{Macro-F1} & \textbf{MASD}$\downarrow$ & \textbf{ECE}$\downarrow$ & \textbf{Repro. recall ($\sigma$)} \\
\midrule
Majority (tillering) & 0.335 & 0.084 & 1.10 & n/a & 0.000 (n/a) \\
ResNet-18, cross-entropy & 0.800 $\pm$ 0.007 & 0.794 $\pm$ 0.019 & 0.238 $\pm$ 0.013 & \textbf{0.062 $\pm$ 0.014} & 0.903 (0.052) \\
ResNet-18, ordinal & \textbf{0.801 $\pm$ 0.013} & \textbf{0.796 $\pm$ 0.018} & \textbf{0.233 $\pm$ 0.013} & 0.064 $\pm$ 0.002 & \textbf{0.917 (0.000)} \\
ResNet-18, focal & 0.767 $\pm$ 0.016 & 0.757 $\pm$ 0.014 & 0.271 $\pm$ 0.010 & 0.179 $\pm$ 0.027 & 0.875 (0.059) \\
\bottomrule
\end{tabular}
\end{table}


% Condensed for 6-page IEEE draft: tab:perclass retained in full manuscript / supplement.


The one place the ordinal loss shows a clear, robust effect is the reproductive stage: its recall is \emph{identical} across all three seeds (0.917, $\sigma=0.000$), whereas cross-entropy swings from 0.833 to 0.958 ($\sigma=0.052$). We therefore frame the ordinal loss's benefit as \emph{stabilizing} reproductive-stage recall rather than increasing it. Per-class recall  shows that this stability does not extend to the other rare class: grain filling, with only 10 test images, is the weakest and most variable class under every loss (recall $0.60$--$0.63$, $\sigma$ up to $0.094$), and the ordinal loss does not improve it. The two under-represented classes therefore behave differently, and neither is claimed to be solved. The focal loss is the clear negative case: worst on accuracy, MASD, and especially calibration (ECE 0.179, roughly three times cross-entropy's), and is kept only as an ablation.

Paired tests on the shared 263-image test split (seed-42 checkpoints) find no significant difference in overall accuracy or MASD between losses. On the reproductive subset ($n=24$), the ordinal model raises recall from 0.833 to 0.917 on this seed (difference $+0.083$, bootstrap 95\% CI $[0.000,0.208]$; McNemar $b{=}2$, $c{=}0$, $p{=}0.50$). Across the three training runs on the \emph{same} 24-image reproductive subset, the ordinal objective produced identical recall (0.917 on every seed); because these repeated predictions reuse the same images they are not independent observations, and we therefore report them descriptively as across-seed \emph{stability} rather than pooling them for a single significance test. The honest reading is that the sample is too small to establish a \emph{higher} reproductive recall; the defensible claim is the zero-variance stability. Because the split is fixed across seeds, the spread in Table~\ref{tab:results} reflects initialization variance alone, resolving the split/seed confound of earlier pilots. One caveat remains: all losses share one cross-entropy-oriented hyperparameter recipe, so the focal result should be read under that recipe.


% Condensed for 6-page IEEE draft: tab:paired retained in full manuscript / supplement.


\subsection{Selective Prediction}

This analysis uses the primary ResNet-18 cross-entropy checkpoint (seed 42), with the temperature fit on the validation split only and never on test. Temperature scaling (fitted $T=0.875$) leaves the test calibration error essentially unchanged (ECE $0.061\to0.061$; Brier improves marginally, $0.301\to0.298$), so we do not present it as a calibration improvement. Its value is realized through selective prediction: allowing the model to abstain on its least-confident cases (Fig.~\ref{fig:riskcoverage}). At full coverage the error rate is 0.205; abstaining on the least-confident ${\sim}46\%$ of images (coverage 0.54) roughly halves it to 0.099, and the normalized area under the risk--coverage curve is 0.126. This motivates a low-confidence uncertainty notice rather than an overconfident recommendation.

\begin{figure}[t]
\centering
\includegraphics[width=0.95\columnwidth]{figures/selective_prediction.png}
\caption{Selective prediction: selective risk ($1-\text{accuracy}$) versus coverage after temperature scaling. Abstaining on the least-confident cases roughly halves the error rate (AURC $=0.126$).}
\label{fig:riskcoverage}
\end{figure}

\subsection{Stage-Gated Retrieval (RQ2)}

Table~\ref{tab:retrieval} reports the aggregate over all 420 auto-derived scenarios, and Table~\ref{tab:retrsig} reports the paired, image-clustered bootstrap differences (P minus baseline) with Holm-adjusted $p$-values. Adaptive soft gating (P) significantly outperforms the ungated baseline (B0) and the fixed-soft baseline (B3) on every metric (all 95\% CIs exclude zero, $p_{\text{Holm}}<0.001$). Against the strongest baseline, query expansion (B1), P is statistically tied on nDCG ($-0.021$, CI $[-0.065,+0.021]$) and recall ($+0.006$, CI $[-0.029,+0.042]$) and loses raw precision ($-0.070$, $p_{\text{Holm}}=0.005$), but reduces SIRR from 0.159 to 0.038---a difference of $-0.121$ ($p_{\text{Holm}}<0.001$), i.e.\ approximately a 75\% relative reduction in stage-incompatible evidence, the property most relevant to safe advisory retrieval. Hard filtering (B2) achieves the lowest SIRR (0.002) but at a substantial nDCG cost. The retrieval results shows the trade-off: P sits near the oracle in the high-nDCG, low-SIRR corner, while B1 attains the best ranking quality at the worst stage safety.

\begin{table}[t]
\centering
\caption{Paired, image-clustered bootstrap differences of the proposed adaptive gating (P) minus each baseline over 420 scenarios, with 95\% CIs and Holm-adjusted $p$-values. The bootstrap resamples whole test images (199 clusters) to respect same-image scenarios. For nDCG/recall/precision, positive favors P; for SIRR (lower is safer), negative favors P.}
\label{tab:retrsig}
\small
\begin{tabular}{llcc}
\toprule
\textbf{Comparison} & \textbf{Metric} & \textbf{Diff (95\% CI)} & \textbf{$p_{\text{Holm}}$} \\
\midrule
P vs.\ B0 & nDCG@5 & $+0.149$ [$+0.122$, $+0.177$] & $<0.001$ \\
P vs.\ B0 & Recall@5 & $+0.080$ [$+0.055$, $+0.105$] & $<0.001$ \\
P vs.\ B0 & SIRR@5 & $-0.098$ [$-0.122$, $-0.074$] & $<0.001$ \\
\addlinespace
P vs.\ B1 & nDCG@5 & $-0.021$ [$-0.065$, $+0.021$] & 0.68 \\
P vs.\ B1 & Recall@5 & $+0.006$ [$-0.029$, $+0.042$] & 0.74 \\
P vs.\ B1 & Precision@5 & $-0.070$ [$-0.111$, $-0.027$] & 0.005 \\
P vs.\ B1 & SIRR@5 & $-0.121$ [$-0.152$, $-0.090$] & $<0.001$ \\
\addlinespace
P vs.\ B3 & nDCG@5 & $+0.049$ [$+0.033$, $+0.066$] & $<0.001$ \\
P vs.\ B3 & Recall@5 & $+0.036$ [$+0.024$, $+0.047$] & $<0.001$ \\
P vs.\ B3 & SIRR@5 & $-0.035$ [$-0.045$, $-0.026$] & $<0.001$ \\
\bottomrule
\end{tabular}
\end{table}

\begin{table}[t]
\centering
\caption{Retrieval over 420 auto-derived scenarios (end-to-end image-based; oracle uses ground-truth stage as an upper bound). Labels are auto-derived from stage-compatibility, not domain-reviewed. Best non-oracle value per column in bold; SIRR lower is better.}
\label{tab:retrieval}
\small
\begin{tabular}{lcccc}
\toprule
\textbf{Method} & \textbf{P@5} & \textbf{R@5} & \textbf{nDCG@5} & \textbf{SIRR@5}$\downarrow$ \\
\midrule
B0: Ungated hybrid & 0.225 & 0.118 & 0.165 & 0.136 \\
B1: Query expansion & \textbf{0.370} & 0.192 & \textbf{0.335} & 0.159 \\
B2: Hard filter & 0.266 & 0.180 & 0.259 & \textbf{0.002} \\
B3: Fixed soft re-ranking & 0.260 & 0.163 & 0.265 & 0.073 \\
P: Adaptive re-ranking & 0.300 & \textbf{0.198} & 0.314 & 0.038 \\
\midrule
Oracle (stage upper bound) & 0.327 & 0.222 & 0.349 & 0.022 \\
\bottomrule
\end{tabular}
\end{table}

Table~\ref{tab:strat} stratifies the same scenarios by whether the image-based prediction is correct, off by one adjacent stage, or off by two or more, and---following the review---now includes the strongest baseline, query expansion (B1), in every stratum. The comparison is more nuanced than a clean ``P wins'' story. On correct cases (330 scenarios), P and B1 are effectively tied on nDCG (0.334 vs 0.332), but P reaches it with far lower SIRR (0.034 vs 0.174). On adjacent-error cases (84 scenarios), B1 actually attains the highest nDCG (0.372 vs P's 0.246), again at markedly higher SIRR (0.074 vs 0.033). So P's advantage in the common correct/adjacent regime ($414/420$ cases) is not ``highest ranking quality'' but the best ranking-\emph{versus}-safety trade-off: competitive nDCG at three-to-five times lower stage-incompatible retrieval. On the 6 non-adjacent errors the picture inverts toward hard filtering, which yields the best non-oracle nDCG (0.184) and the lowest SIRR (0.100 versus P's 0.300), because when the vision model is confidently wrong by several stages soft gating overweights the incorrect stage. Although $n=6$ is too small for a statistical claim, it shows adaptive gating is not uniformly safest---its value is concentrated where the belief is right or close.

\begin{table}[t]
\centering
\caption{nDCG@5 / SIRR@5 stratified by vision error condition (420 scenarios), including the strongest baseline B1. On correct/adjacent cases B1 attains the highest raw nDCG but at markedly higher SIRR; adaptive gating (P) offers the best ranking-versus-safety trade-off, while hard filtering (B2) is safest under rare non-adjacent errors. Bold marks the best non-oracle nDCG and the lowest SIRR per stratum.}
\label{tab:strat}
\small
\begin{tabular}{lccc}
\toprule
\textbf{Method} & \textbf{Correct ($n{=}330$)} & \textbf{Adjacent ($n{=}84$)} & \textbf{Non-adj. ($n{=}6$)} \\
\midrule
B0: Ungated & 0.162 / 0.146 & 0.185 / 0.074 & 0.052 / 0.433 \\
B1: Query expansion & 0.332 / 0.174 & \textbf{0.372} / 0.074 & 0.000 / 0.500 \\
B2: Hard filter & 0.276 / \textbf{0.000} & 0.195 / \textbf{0.005} & \textbf{0.184} / \textbf{0.100} \\
B3: Fixed soft & 0.277 / 0.073 & 0.228 / 0.057 & 0.135 / 0.300 \\
P: Adaptive & \textbf{0.334} / 0.034 & 0.246 / 0.033 & 0.174 / 0.300 \\
\bottomrule
\end{tabular}
\end{table}


% Trade-off visualization kept in full manuscript; IEEE draft retains numerical retrieval table.


\section{Discussion}
Under a single fixed group split, a ResNet-18 achieves about 0.80 accuracy against \emph{weakly derived} BBCH macro-stage labels. Ordinal loss changes overall accuracy negligibly but gives identical reproductive recall (0.917) across three seeds on the same 24 test images; this is descriptive within-split stability, not statistically established improvement or cross-field generalization. Grain filling (10 test images) remains underpowered.

Adaptive soft gating offers a ranking-versus-incompatibility operating point, not uniform dominance: against query expansion its reported nDCG is 0.314 versus 0.335 and its SIRR is 0.038 versus 0.159, with lower precision. Hard filtering retains the lowest SIRR; rare non-adjacent visual errors may favor hard filtering. Because relevance and incompatibility are derived from the same stage-compatibility metadata consumed by the ranker, results validate internal stage metadata consistency rather than independently judged agronomic effectiveness.

\section{Limitations and Future Work}
The 2,900 RiceSEG-derived macro-stage labels depend on canopy cover, panicle fraction and senescence thresholds, not independent expert BBCH annotation; 1,134 images are flagged for review. Hash-aware parent grouping prevents detected aHash near-duplicate pairs crossing splits at the configured threshold but does not establish new-field, season or device independence. The fixed-split three-seed protocol measures training variability conditioned on one partition, and post-hoc selective prediction operating points require validation-based threshold selection. Future work is expert BBCH review, independent relevance judgments, acquisition-grouped external testing, calibrated threshold selection and systematic ablations of adaptive stage weight versus a validation-tuned fixed weight. Chemical or regulatory advice remains outside unattended deployment.

\section{Ethical Considerations}

Farm images can reveal location, farm layout, or other information the grower may not want disclosed, so any public release should strip personal identifiers and precise geolocation unless the source has explicitly consented. The system is a research component feeding a larger decision-support pipeline, not an autonomous authority on agronomic decisions. Outputs involving pesticides, dosage, withholding periods, or regulated actions must be grounded in current authoritative sources, and low-confidence predictions should prompt a request for a better image or an uncertainty flag rather than a definitive-sounding recommendation.

\section{Conclusion}

CropState-Vision recognizes six BBCH-derived rice growth stages from field images and uses the full stage distribution to gate evidence retrieval. On a 2{,}900-image benchmark whose splits show no detected cross-split near-duplicates at perceptual-hash thresholds up to 10, a fixed-split, multi-seed study shows a fine-tuned ResNet-18 reaches 0.800 accuracy, and among training objectives an ordinal-aware loss stabilizes reproductive-stage recall without changing accuracy while a focal loss degrades calibration. Casting confidence as selective prediction roughly halves the error rate when the model abstains on its least-confident cases. On a preliminary automated evaluation over 420 auto-derived scenarios, adaptive soft gating stays competitive with the strongest baseline in ranking quality (statistically tied on nDCG and recall under an image-clustered bootstrap) while cutting stage-incompatible retrievals by about 75\%, with a clearly bounded exception under rare non-adjacent errors. Two honesty caveats govern these claims: the reproductive-recall result is across-seed stability rather than a significant increase at this sample size, and the retrieval labels are auto-derived from the same metadata the re-ranker uses, making the study a mechanism validation rather than domain-reviewed effectiveness. An expert-annotated relevance set and field- or season-grouped metadata are the two items that would convert these findings into deployment-grade evidence.

\section{Data and Code Availability}
\label{sec:availability}

The image benchmark is derived from the publicly available RiceSEG dataset; we use it for research under its original terms and do not redistribute the raw images. Our contribution over the source is the six weak macro-stage labels inferred from RiceSEG segmentation-mask morphology, the fixed group split (2{,}400/237/263, seed 42), and the near-duplicate audit. The knowledge corpus is built from the publicly available IRRI Rice Knowledge Bank and is used under the source's terms. To support reproduction, the stage mapping, the split manifest, the derived stage-compatibility tags, the training and evaluation code (including the perceptual-hash audit, the selective-prediction analysis, and the image-clustered bootstrap used for Table~\ref{tab:retrsig}), and the aggregated result files will be released in a public repository upon publication, or are available from the authors on request pending license and metadata clean-up. No expert-annotated relevance set is released because none yet exists; constructing one is stated future work.

\bibliographystyle{IEEEtran}
\bibliography{references}

\end{document}
